\subsection{Completion and whole-figure supports}
Let $E$ be the event poset obtained in Theorem~\ref{thm:events}.
Thus conflict $\#$ is symmetric, irreflexive and hereditary upward,
and every principal past $\downarrow e$ is finite. A configuration
is a downward-closed conflict-free subset. Write $\mathcal C_f(E)$
for finite configurations and $\mathcal C(E)$ for all configurations.
The orders are inclusion, with empty configuration $0$.

For two configurations define whole-figure disjointness by
\[
 F\perp G\quad\Longleftrightarrow\quad
 F\cap G=\varnothing\ \text{ and }\ F\cup G\in\mathcal C(E).
\]
On events use the relation
\[
 e\perp_E f\quad\Longleftrightarrow\quad
 \neg(e\# f)\ \text{ and }\
 \downarrow e\cap\downarrow f=\varnothing.
\]
These relations concern whole figures. They are different from the
primitive disjunction between open cells used in Section~7.

\begin{proposition}[Configuration completion and supports]
\label{prop:event-domain}
There is a canonical order isomorphism
\[
 \mathcal C(E)\simeq\operatorname{Idl}(\mathcal C_f(E)),
\]
where an ideal is a nonempty directed downward-closed subset.
The compact elements of $\mathcal C(E)$ are its finite configurations,
and its nonzero completely join-prime elements are exactly the
principal pasts $\downarrow e$. Every pairwise bounded family has
a join. Moreover there are canonical complete ortholattice
isomorphisms
\[
 \Sigma(E,\perp_E)\simeq
 \Sigma(\mathcal C_f(E)\setminus\{0\},\perp)
 \simeq\Sigma(\mathcal C(E)\setminus\{0\},\perp).
\]
If $0$ is retained as a universally self-orthogonal vertex,
the same support lattice results, with bottom support $\{0\}$.
\end{proposition}
The configuration-domain assertions are the specialization of
Nielsen--Plotkin--Winskel's coherent prime-algebraic domain
correspondence \cite[Definitions~17--18 and Theorems~8--9, p.~102]{NPW}.
We give the elementary argument and the support comparison with
all empty and infinite cases explicit.

\begin{proof}
Every principal past is conflict-free: a conflict between two of
its elements would propagate to a self-conflict at its maximum.
A family of configurations is bounded exactly when its union is
conflict-free. The union is then its join. Pairwise boundedness
already excludes every possible conflict, proving the coherence
assertion and the formula for directed joins.

For a configuration $C$, its finite subconfigurations form an ideal.
Their union is $C$, since each event belongs to its finite principal
past. Conversely the union of an ideal is a configuration: any
two events lie together in a common ideal member. A finite
subconfiguration of that union lies in one ideal member by
directedness and therefore belongs to the ideal. These operations
are inverse, including when $E$ is empty.

A finite configuration contained in a directed union is contained
in one member, by finitely many witnesses and directedness. It is
therefore compact. Every configuration is the directed union of
the finite unions of its principal pasts, so a compact configuration
must itself be finite. If $\downarrow e\subseteq\bigvee_i C_i$,
the join is a union and $e\in C_i$ for some $i$. Downward closure
gives $\downarrow e\subseteq C_i$. Conversely a nonzero completely
join-prime configuration, written as the union of its principal
pasts, must equal one of them. This proves the domain assertions.

The event orthogonality is hereditary downward in both variables.
For configurations, finite or infinite, one has
\begin{equation}\label{eq:event-pairwise}
 F\perp G\quad\Longleftrightarrow\quad
 (\forall e\in F)(\forall f\in G)\ e\perp_E f.
\end{equation}
In the forward direction causal closure puts the two pasts in
the two disjoint configurations. Conversely the displayed condition
excludes cross-conflicts and any common event, proving the claim.

Let $V$ be either nonempty configuration set in the statement.
For $\mathcal U\subseteq V$ and $S\subseteq E$, put
\[
 U(\mathcal U)=\bigcup_{F\in\mathcal U}F,
 \qquad J(S)=\{F\in V:F\subseteq S\}.
\]
Equation~\eqref{eq:event-pairwise} gives
\[
 N_V(\mathcal U)=J\bigl(N_E(U(\mathcal U))\bigr).
\]
Every event polar is downward closed. If $S$ is downward closed,
then $U(J(S))=S$, because the finite past $\downarrow e$
belongs to $J(S)$ for every $e\in S$. Consequently, for a closed
event support $S$,
\[
 N_V(J(S))=J(N_E(S)),\qquad N_V^2(J(S))=J(S).
\]
Conversely, applying the same formulas twice to a closed
configuration support $\mathcal U$ yields
$\mathcal U=J(N_E^2(U(\mathcal U)))$ and
$U(\mathcal U)=N_E^2(U(\mathcal U))$.
Thus $J$ and $U$ are inverse order and orthocomplement isomorphisms.
This proves the assertion for both finite and arbitrary configurations.
The test objects used here are principal pasts; a singleton event
need not be a configuration.

Finally, if $0$ is retained, every polar contains it because it
is orthogonal to every figure, including itself. Adding or removing
this one vertex gives the asserted support isomorphism.
\end{proof}

\subsection{Which data a configuration domain forgets}
Incidence and compatibility are recovered from configurations.
Refinement requires additional structure. For clarity the
translation of the refinement-by-strata rule is
\[
 G\ll F\quad\Longleftrightarrow\quad
 (\forall e\in G)(\exists f\in F)\ e\ll f,
\]
where refinement of events means refinement of their principal
configurations. If subdivision is defined to preserve whole-figure
cosupport, its translation is
\[
 G\preceq F\quad\Longleftrightarrow\quad
 G\ll F\ \text{ and }\ N(G)=N(F).
\]
Here $N(F)$ denotes all configurations disjoint from $F$.
Triple-polar equality makes this equivalent to equality of the
double-polar supports. Neither formula supplies the additional
refinement relation from the event structure alone.

The following examples use a precise elementary convention. The
principal deployment of a stratum $f$ is $\downarrow f$. A carrier
of $e\ll f$ is the least $d\leq f$ with $e\ll d$; internal
refinement $e\rho f$ means that this carrier is $f$ itself.
A primitive disjunction is symmetric and irreflexive, is hereditary
under internal refinement in both variables, and relates every
two distinct members of one principal deployment. Figures in the
examples are exactly its downward-closed cliques.

For an explicit example, take two incomparable conflicting events
$a,b$. The configurations are $0,A=\{a\},B=\{b\}$.
Use incidence alone as refinement in the first system, and add
$A\ll B$ in the second. Both are partial orders; their nonzero
principal deployments are singletons, so carriers exist uniquely,
and compatible refinement is incidence. Internal refinement is
transitive. Set primitive disjunction on the two strata to be empty;
its heredity and the distinct-strata condition within one principal
deployment are vacuous. The first system has two minimal refinement
objects, or lieux, and the second has only $a$. All nonzero figures
have polar $\{0\}$, so $A$ subdivides $B$ only in the second system.

Primitive internal disjunction can also vary while the event
structure stays fixed. Take two chains $a<a'$ and $b<b'$, with
every cross-chain pair in conflict. The configurations are
\[
 0,\quad\{a\},\quad\{a,a'\},\quad\{b\},\quad\{b,b'\}.
\]
Set refinement equal to incidence. One primitive relation has
exactly the unordered pairs $\{a,a'\}$ and $\{b,b'\}$; the other
also has $\{a',b'\}$. Internal refinement is identity, so heredity
holds. In each case the downward-closed primitive cliques are
exactly the five displayed configurations: any mixed downset
contains the missing pair $\{a,b\}$. Thus the causal order,
conflict, configurations and whole-figure supports agree, while
the primitive relation differs.

This also distinguishes a cosupport definition from a maximality
theorem. A witness-absorption condition saying that a figure
disjoint from $P\ll F$ but not from $F$ can be chosen to refine
$F$ does not alone make subdivisions incidence-maximal. In the
last example it holds: for nonzero $P$ its premise is empty, and
for $P=0$ one can use $F$ itself. Yet $\{a\}$ is a proper
cosupport-preserving incidence subfigure of $\{a,a'\}$.
The upper strata in this example have no open lieux. Under the
strong carrier hypotheses, nonempty open-lieu sets and witness
absorption give the exact maximality conclusion of
Theorem~\ref{thm:abstract-extension}; the full atlas then also
gives subdivision extension.
